\chapter{Analyse et spécifications}
\label{chap:analyse}

\section{Périmètre}

L'application cible est un \emph{screensaver} interactif présentant
un aquarium peuplé de plusieurs espèces de poissons. Elle s'exécute
sur un poste Linux avec un environnement graphique compatible GTK4.
Elle est mono-utilisateur, mono-fenêtre, et n'a pas de persistance~:
chaque démarrage recharge la configuration depuis le disque et
réinitialise la simulation.

\section{Exigences fonctionnelles}

Les exigences fonctionnelles, numérotées EF1..EF12, décrivent ce
que le système doit faire.

\begin{description}[leftmargin=3em, style=nextline]
  \item[EF1 -- Simulation de banc.] L'application doit simuler le
    comportement de banc d'au moins sept espèces de poissons
    distinctes, en appliquant à chacune les trois forces de
    Reynolds (séparation, alignement, cohésion) avec des
    paramètres propres.

  \item[EF2 -- Prédateur/proie.] Certaines espèces sont des
    prédateurs et chassent d'autres espèces désignées comme leurs
    proies. La consommation d'une proie est effective lorsque la
    \emph{bouche} du prédateur entre en collision avec le centre
    de la proie.

  \item[EF3 -- Fuite des proies.] Toute proie qui détecte un
    prédateur dans son rayon de peur (\code{fear\_radius}) doit
    s'enfuir en sens opposé, avec un \emph{boost} de vitesse
    d'adrénaline.

  \item[EF4 -- Cycle de vie.] Une entité passe par les états
    \code{ALIVE}, \code{DYING} (transition visuelle de
    $\sim$~0.22~seconde) puis \code{INACTIVE} (invisible,
    en attente). Une entité \code{INACTIVE} peut être réactivée
    au-delà du délai \code{respawn\_delay}.

  \item[EF5 -- Famine.] Un prédateur qui n'a rien mangé depuis
    \code{STARVATION\_SECONDS} secondes meurt de faim.

  \item[EF6 -- Configuration externe.] L'intégralité des
    paramètres (constantes globales et tableau des espèces) est
    lue au démarrage depuis un fichier TOML
    \file{aquarium.conf}.

  \item[EF7 -- Recherche multi-niveaux.] Le fichier de
    configuration est cherché successivement dans le répertoire
    courant, puis dans \file{\textasciitilde/.config/gtkaqua/},
    puis dans \file{/etc/gtkaqua/}.

  \item[EF8 -- Ajout/suppression d'espèces.] Ajouter une espèce
    consiste à ajouter un bloc \code{[[species]]} au fichier de
    configuration et à fournir le sprite correspondant dans le
    répertoire \file{assets/}, sans modifier le code C ni
    recompiler.

  \item[EF9 -- Validation.] Le chargement doit échouer proprement
    en cas d'erreur de syntaxe TOML, d'identifiant manquant, ou
    de référence \code{prey} pointant vers une espèce inexistante,
    en affichant un message lisible.

  \item[EF10 -- Rechargement à chaud.] Un bouton \emph{Reload
    config} dans l'interface relit le fichier de configuration et
    relance la simulation avec les nouveaux paramètres, sans
    redémarrer l'application.

  \item[EF11 -- Contrôles utilisateur.] L'interface latérale
    fournit~: une case à cocher par espèce pour décider qui peuple
    le \emph{pool} initial, un sélecteur d'espèce et un compteur
    pour ajouter/retirer un nombre précis de poissons à chaud, un
    bouton pause, un curseur de vitesse globale, un bouton
    \emph{Initialize pool}, le bouton \emph{Reload config}.

  \item[EF12 -- Plein écran.] La touche \kw{F11} bascule en mode
    plein écran. La touche \kw{Échap} quitte le plein écran.
\end{description}

\section{Exigences non fonctionnelles}

\begin{description}[leftmargin=3em, style=nextline]
  \item[ENF1 -- Performance.] La simulation doit tenir 60~Hz avec
    au moins 100 entités sur un poste de bureau moderne (CPU
    quad-core, GPU intégré).

  \item[ENF2 -- Portabilité.] Cible Linux x86\_64 avec GTK4 et
    GLib installés. Pas d'objectif de portage Windows ou macOS.

  \item[ENF3 -- Lisibilité de la configuration.] Le fichier
    \file{aquarium.conf} doit être éditable par un utilisateur
    qui n'est pas développeur~: noms de paramètres explicites,
    commentaires de section, valeurs typées (entiers, flottants,
    listes de chaînes), pas de syntaxe spécifique au langage~C.

  \item[ENF4 -- Robustesse mémoire.] Pas de fuite de mémoire en
    régime stationnaire. Les widgets GTK des entités
    \code{INACTIVE} sont conservés et masqués plutôt que détruits
    et recréés, pour éviter la fragmentation et les coûts
    d'allocation répétés.

  \item[ENF5 -- Modularité.] Chaque module C correspond à un
    couple \file{.h}/\file{.c} clairement délimité, et expose
    une API minimale. La taille moyenne d'un module reste
    raisonnable (quelques centaines de lignes).

  \item[ENF6 -- Reproductibilité du build.] La compilation se
    fait par \code{make} sans télécharger de dépendances
    réseau au moment du build. Le seul \emph{vendor} embarqué
    est la bibliothèque \code{tomlc99}, copiée dans
    \file{vendor/tomlc99/}.

  \item[ENF7 -- Diagnostic.] En cas d'erreur de chargement de
    configuration, un message clair indique le chemin essayé et
    la cause (clé manquante, identifiant de proie inconnu, etc.).
\end{description}

\section{Cas d'utilisation}

Cinq cas d'utilisation principaux structurent l'interaction avec
l'application. Ils sont représentés graphiquement dans la
figure~\ref{fig:usecases}.

\begin{itemize}[leftmargin=*]
  \item \textbf{UC1 -- Lancer la simulation.} L'utilisateur
    démarre l'application. Le programme lit
    \file{aquarium.conf}, instancie les espèces, peuple le
    \emph{pool} initial selon \code{flock\_size} de chaque
    espèce, démarre la boucle de simulation.

  \item \textbf{UC2 -- Mettre en pause.} L'utilisateur clique sur
    \emph{Pause}. Le tick continue d'être appelé mais
    \code{world\_tick} n'est plus exécuté.

  \item \textbf{UC3 -- Modifier l'écosystème via configuration.}
    L'utilisateur édite \file{aquarium.conf} (modification d'un
    paramètre, ajout/suppression d'un bloc \code{[[species]]}),
    enregistre, puis clique sur \emph{Reload config} ou redémarre
    l'application.

  \item \textbf{UC4 -- Ajouter/retirer des poissons à chaud.}
    L'utilisateur sélectionne une espèce dans la combobox,
    saisit un nombre, et clique sur \emph{+~Add} ou
    \emph{-~Remove}.

  \item \textbf{UC5 -- Recharger la configuration.} L'utilisateur
    clique sur \emph{Reload config}. Le système relit le fichier,
    valide, et relance la simulation. En cas d'erreur, une boîte
    de dialogue s'affiche.
\end{itemize}

\begin{figure}[h!]
\centering
\begin{tikzpicture}[
  actor/.style = {draw, thick, circle, minimum size=0.7cm, label=below:{\small Utilisateur}},
  ucase/.style = {draw, thick, ellipse, minimum width=3.6cm, minimum height=1cm, align=center, font=\small},
  sysbox/.style = {draw, thick, rounded corners, inner sep=0.6cm, label=above:{\small \textbf{Aquarium GTK}}}
]

\node[actor] (user) at (0,-2.5) {};

\node[ucase] (uc1) at (5, 0)    {UC1 : Lancer\\la simulation};
\node[ucase] (uc2) at (5, -1.5) {UC2 : Mettre\\en pause};
\node[ucase] (uc3) at (5, -3)   {UC3 : Modifier\\la configuration};
\node[ucase] (uc4) at (5, -4.5) {UC4 : Ajouter/retirer\\des poissons};
\node[ucase] (uc5) at (5, -6)   {UC5 : Recharger\\la configuration};

\node[sysbox, fit=(uc1)(uc2)(uc3)(uc4)(uc5)] {};

\foreach \uc in {uc1, uc2, uc3, uc4, uc5} {
  \draw[thick] (user) -- (\uc);
}

\end{tikzpicture}
\caption{Diagramme de cas d'utilisation.}
\label{fig:usecases}
\end{figure}

\section{Glossaire}

\begin{description}[leftmargin=3em, style=nextline]
  \item[Boid.] Agent autonome dans le modèle de Reynolds. Dans
    notre projet, un boid est une instance de \struct{Entity}.

  \item[Flocking.] Comportement de banc émergeant de
    l'application des trois règles de Reynolds.

  \item[Mouth-based eating.] Approche d'alimentation où la
    collision est calculée entre la \emph{bouche} du prédateur
    (point décalé vers l'avant du sprite, à
    \code{mouth\_forward\_ratio}~$\times$~\code{sprite\_width}) et
    la position centrale de la proie. Plus réaliste qu'une
    collision centre-à-centre.

  \item[Lifespan.] Durée de vie maximale d'une entité, exprimée
    en secondes. Atteinte la durée, l'entité passe à
    \code{DYING}.

  \item[Respawn.] Réactivation d'une entité \code{INACTIVE} après
    expiration de \code{respawn\_delay} et tant que la population
    vivante de l'espèce est inférieure à \code{flock\_size}.

  \item[Pool.] Ensemble des espèces effectivement peuplées au
    démarrage, déterminé par les cases à cocher de l'interface.

  \item[Pred mask.] Bitmask 32~bits où le bit~$i$ indique que
    l'espèce d'index~$i$ est une proie de l'espèce courante.
\end{description}
